\documentclass[../../../include/open-logic-section]{subfiles}

\begin{document}

\olfileid{sth}{story}{rus}
\olsection{رسل دا پیراڈوکس (دوبارہ)}

\olref[sfr][][]{part} وچ اسیں اک سادہ سیٹ نظریے نال کم کیتا سی۔ پر اک
\emph{بہت ای} سادے خیال مطابق سیٹ صرف محمولاں دے دائرے ہوندے نیں۔ ایہ سادہ خیال ہِیٹھلا اصول منّن
لئی آکھدا اے:

\begin{defish}
  \emph{سادہ خاصیت توں سیٹ بناؤن دا اصول۔} $\Setabs{x}{\phi(x)}$ ہر فارمولے $\phi$ لئی موجود اے۔
\end{defish}

بھانویں ایہ اصول منّن نوں دل کردا اے، پر ثابت کیتا جا سکدا اے کہ ایس
وچ تضاد اے۔ ایہ اسیں \olref[sfr][set][rus]{sec} وچ ویکھیا سی۔ پر نتیجہ اینا
اہم تے اینا سدھا اے کہ اینوں دوبارہ کہنا بندا اے۔ عین اوہناں لفظاں وچ۔

\begin{thm}[رسل دا پیراڈوکس]
کوئی سیٹ $R = \Setabs{x}{x \notin x}$ موجود نہیں۔
\end{thm}

\begin{proof}
جے $R = \Setabs{x}{x \notin x}$ موجود ہووے، تاں
$R \in R$ اُدوں تے صرف اُدوں اے جدوں $R \notin R$، تے ایہ تضاد اے۔
\end{proof}

رسل نے ایہ نتیجہ جون 1901 وچ لبھیا۔ (پر اوہنے پیراڈوکس عین اوس صورت وچ
پیش نہیں کیتا جیویں اسیں ہُنے کیتا اے۔ اوہ فریگے دے سیٹ نظریے بارے سوچ ریا سی،
جیہڑا \emph{Grundgesetze} وچ بیان ہویا اے۔ اسیں \olref[blv]{sec} وچ ایس گل
ول مُڑاں گے۔) رسل نے 16 جون 1902 نوں فریگے نوں خط لکھ کے اوہدے نظام دا
تضاد سمجھایا۔ خط کتابت تے کجھ پچھوکڑ لئی \citet[pp.~124--8]{Heijenoort1967} ویکھو۔

ایہ گل زور دے کے آکھنی چاہیدی اے کہ دو سطراں والا ایہ ثبوت \emph{خالص منطق}
دا نتیجہ اے۔ منّیا کہ اسیں اپنی علامت $\Setabs{x}{x \notin x}$ وچ اک
(غیر منطقی؟) مسلّمہ، عنصراں نال تعیّن، لُکواں ورتیا اے۔ کیوں جے
$\Setabs{x}{\phi(x)}$ دا مطلب اے کہ جے کوئی سیٹ موجود اے تاں اوہ $\phi$
والیاں چیزاں دا \emph{اکّو اک} سیٹ اے (عنصراں نال تعیّن دے اصول مطابق)۔
پر عنصراں نال تعیّن دا اشارہ وی ہٹ سکدا اے جے نتیجہ ایویں آکھیے:
\emph{کوئی ایہو جیہا سیٹ نہیں جیہدے رکن عین اوہ سیٹ ہون جیہڑے اپنے آپ دے رکن نہیں}۔
تے ایہ گل خاص طور اُتے سیٹاں نال جڑی ہوئی نہیں۔ جیویں رسل نے آپ دھیان دوایا،
عین ایہو جیہی دلیل نال ایہ نتیجہ نکلدا اے کہ \emph{کوئی مرد عین اوہناں
مرداں دی داڑھی نہیں مونڈدا جیہڑے اپنی داڑھی آپ نہیں مونڈدے}۔ یا
\emph{کوئی پگ کُتّا عین اوہناں پگ کُتیاں نوں نہیں سُنگھدا جیہڑے اپنے آپ نوں
نہیں سُنگھدے}۔ تے ایویں ہور مثالاں۔ نتیجے دی عام صورت بس ایہ اے:
\[
\lnot \exists x \forall z(Rzx \liff \lnot R zz).
\]
تے ایہ پہلے درجے دی منطق دا اک قضیہ (سانچہ) اے۔ ایس لئی سیٹ نظریے وچ
تھوڑی جیہی ردّوبدل کر کے رسل دے پیراڈوکس توں نہیں بچ سکدے۔ ایہ سیٹ نظریے
تک \emph{پہنچن} توں پہلاں ای پیدا ہوندا اے۔ جے اسیں (کلاسیکی) پہلے درجے دی
منطق ورتنی اے، تاں سانوں ایہ \emph{منّنا} پئے گا کہ کوئی سیٹ $R = \Setabs{x}{x\notin x}$ موجود نہیں۔

نتیجہ ایہ اے۔ جے تسی سادہ خاصیت توں سیٹ بناؤن دا اصول منّنا چاہندے او
تے تضاد توں \emph{بچنا} وی، تاں تسی صرف \emph{سیٹ نظریے} وچ ردّوبدل نال
کم نہیں چلا سکدے۔ تہانوں اپنی \emph{منطق} دی بناوٹ بدلنی پئے گی۔

بے شک غیر کلاسیکی منطقاں والے سیٹ نظریے پیش کیتے گئے نیں۔ پر ایہ
گھٹ توں گھٹ اینا تاں آکھ سکدے آں کہ اوہ غیر معیاری نیں۔ رسل دے پیراڈوکس
نال نِبڑن دا معیاری طریقہ اینوں موجود نہ ہون دا سدھا ثبوت منّنا تے فیر
ایہدے نال کم کرنا سِکھنا اے۔ اسیں ایہو طریقہ ورتاں گے۔

\end{document}
